% GENERATED FILE. Edit canonical lessons, then run npm run generate:books.
% canonical-source-hash: fnv1a64:9e7f59f6669daf9c
% canonical-lessons: ZH-C173-yaoqing, ZH-C173-wenhou, ZH-W173-hou-wait, ZH-C173-wanhui, ZH-C173-juhui, ZH-C173-qingke

\chapter{An Invitation}
\label{ch:zh-an-invitation}
\hlchaptermodality{173}

\begin{chapteropening}
\textbf{By the end of this chapter:} \emph{I can invite people, send my regards, and talk about an evening party, a get-together and who pays.}

Complete the last lesson of chapter 173: I can invite people, send my regards, and talk about an evening party, a get-together and who pays.
\end{chapteropening}

\section[yāoqǐng]{yāoqǐng — to invite; an invitation}

\label{lesson:ZH-C173-yaoqing}



\begin{quote}
Before the new one: say the Mandarin for smelling good, then the Mandarin for bland.
\end{quote}

\subsection*{You'll want to know: yāoqǐng}
\emph{yāoqǐng} — \textquotedblleft{}to invite; an invitation\textquotedblright{}.

\begin{quote}
\emph{tā yāoqǐng wǒmen qù tā jiā} — He's invited us to his place.
\end{quote}

Its characters come later; for now it is a word for the ear.

\begin{cousinweb}
\emph{yāoqǐng} is built from \emph{yāo} and \zh{请} 'to ask', both inherited from Middle Chinese, \emph{'jiew} and \emph{tshjengX} (Guangyun, in Baxter's transcription); the dictionary glosses the first 'to stop someone on the way'.
\end{cousinweb}

\begin{grammarlens}[title={What you've built}]
One more word for talking about plans.
\end{grammarlens}

\subsection*{Your turn}
\begin{itemize}
\raggedright
  \item \emph{Say it:} \emph{yāoqǐng}
  \item \emph{Say it:} \emph{yāoqǐng}, once more
  \item \emph{Say it:} \emph{tā yāoqǐng wǒmen qù tā jiā} — He's invited us to his place.
  \item \emph{Recall:} say the Mandarin for smelling good, then the Mandarin for bland, then say \emph{yāoqǐng} again
\end{itemize}

\begin{tcolorbox}[breakable,colback=teal!4,colframe=teal!35!black,title={Before you move on}]
What does yāoqǐng mean? (\textquotedblleft{}To invite; an invitation\textquotedblright{}.) Say it once more.
\end{tcolorbox}
\section[\texorpdfstring{\zh{问候} (wènhòu)}{问候 (wènhòu)}]{\zh{问候} (wènhòu) — to send regards to; to greet}

\label{lesson:ZH-C173-wenhou}



\begin{quote}
Before the new one: say the Mandarin for bland, then the Mandarin for to invite.
\end{quote}

\subsection*{You'll want to know: \zh{问候}}
\textbf{\zh{问候}} — \emph{wènhòu} — \textquotedblleft{}to send regards to; to greet\textquotedblright{}.

\begin{quote}
\emph{nǐ bāng wǒ wènhòu nǐ bàba māma} — Give my regards to your mum and dad.
\end{quote}

The new character, in pieces: \zh{候} — ten strokes: the person piece, a thin vertical beside it, then on the right a small hooked top, a horizontal, and a base of five strokes: a short stroke falling to the left, two horizontals, and two long strokes falling to either side. The next lesson puts it in your hand.

\begin{cousinweb}
\textbf{\zh{问候}} (\emph{wènhòu}) is built from \zh{问} 'to ask' and \zh{候} 'to wait upon', both inherited from Middle Chinese, \emph{mjunH} and \emph{huwH} (Guangyun, in Baxter's transcription): to ask after someone. The second is a different word from \zh{后} 'after', \emph{huwX}.
\end{cousinweb}

\begin{grammarlens}[title={What you've built}]
One more word for talking about plans.
\end{grammarlens}

\subsection*{Your turn}
\begin{itemize}
\raggedright
  \item \emph{Say it:} \emph{wènhòu}
  \item \emph{Say it:} \emph{wènhòu}, once more
  \item \emph{Say it:} \emph{nǐ bāng wǒ wènhòu nǐ bàba māma} — Give my regards to your mum and dad.
  \item \emph{Recall:} say the Mandarin for bland, then the Mandarin for to invite, then say \emph{wènhòu} again
\end{itemize}

\begin{tcolorbox}[breakable,colback=teal!4,colframe=teal!35!black,title={Before you move on}]
What does \zh{问候} mean? (\textquotedblleft{}To send regards to; to greet\textquotedblright{}.) Say it once more.
\end{tcolorbox}
\section[\texorpdfstring{\zh{候} (hòu)}{候 (hòu)}]{\zh{候} — the person piece, a thin stroke, then the right side}

\label{lesson:ZH-W173-hou-wait}



\begin{quote}
The word first, then the character inside it.

\begin{itemize}
\raggedright
  \item \emph{Recall:} say \emph{wènhòu} — \textbf{R1}, one lesson back, before it decays
\end{itemize}

\begin{quote}
\zh{亻} \zh{问}
\end{quote}

\emph{Read it:} each one aloud, with nothing said first — the person piece, \emph{wèn}
\end{quote}

\subsection*{Script — the person piece, a thin stroke, then the right side}
\hlblockfigure{\detokenize{figures/ZH-W173-hou-wait-filmstrip.pdf}}{How \zh{候} is written, stroke by stroke}

\begin{quote}
\zh{候}
\end{quote}

\textbf{Three pieces, left to right.} Count the pieces, not the strokes.

\textbf{The person piece, \zh{亻}, two strokes,} as in \zh{你}.

\textbf{A thin vertical,} one short stroke standing beside it.

\textbf{The right side, seven strokes.} On top a short horizontal that turns and falls to the left, like a small hook. A horizontal under it. Then the base, five strokes: a short stroke falling to the left, a short horizontal, a longer horizontal, a long stroke falling to the left through them, and a long stroke falling to the right.

2 + 1 + 7 = ten.

\subsection*{Your turn}
\begin{itemize}
\raggedright
  \item \emph{Write it:} \zh{候}, all ten, with the model in front of you
  \item \emph{Cover:} write it again from memory, and count the pieces, not the strokes
\end{itemize}

\begin{tcolorbox}[breakable,colback=teal!4,colframe=teal!35!black,title={Before you move on}]
How many strokes in \zh{候}? (\textbf{Ten}.) What comes after the person piece? (\textbf{One thin vertical}.) How many strokes in the base? (\textbf{Five}.)

Source: \href{https://raw.githubusercontent.com/chanind/hanzi-writer-data/68d10a4b21150cae5e1ebbd223eed289cf32d90c/data/\%E5\%80\%99.json}{Hanzi Writer Data \zh{候}, pinned stroke paths and medians}.
\end{tcolorbox}
\section[\texorpdfstring{\zh{晚会} (wǎnhuì)}{晚会 (wǎnhuì)}]{\zh{晚会} (wǎnhuì) — evening party; evening show}

\label{lesson:ZH-C173-wanhui}



\begin{quote}
Before the new one: say the Mandarin for to invite, then the Mandarin for to send regards to.
\end{quote}

\subsection*{You'll want to know: \zh{晚会}}
\textbf{\zh{晚会}} — \emph{wǎnhuì} — \textquotedblleft{}evening party; evening show\textquotedblright{}.

\begin{quote}
\textbf{\zh{晚上我有事不能去晚会}} \emph{wǎnshang wǒ yǒushì bù néng qù wǎnhuì} — I've got something on tonight, so I can't come to the party.
\end{quote}

\begin{cousinweb}
\textbf{\zh{晚会}} (\emph{wǎnhuì}) is built from \zh{晚} 'evening' and \zh{会} 'to meet', both inherited from Middle Chinese, \emph{mjonX} and \emph{hwajH} (Guangyun, in Baxter's transcription), which glosses the second 'to come together'.
\end{cousinweb}

\begin{grammarlens}[title={What you've built}]
One more word for talking about plans.
\end{grammarlens}

\subsection*{Your turn}
\begin{itemize}
\raggedright
  \item \emph{Say it:} \emph{wǎnhuì}
  \item \emph{Say it:} \emph{wǎnhuì}, once more
  \item \emph{Say it:} \emph{wǎnshang wǒ yǒushì bù néng qù wǎnhuì} — I've got something on tonight, so I can't come to the party.
  \item \emph{Recall:} say the Mandarin for to invite, then the Mandarin for to send regards to, then say \emph{wǎnhuì} again
\end{itemize}

\begin{tcolorbox}[breakable,colback=teal!4,colframe=teal!35!black,title={Before you move on}]
What does \zh{晚会} mean? (\textquotedblleft{}Evening party; evening show\textquotedblright{}.) Say it once more.
\end{tcolorbox}
\section[jùhuì]{jùhuì — get-together; party}

\label{lesson:ZH-C173-juhui}



\begin{quote}
Before the new one: say the Mandarin for to send regards to, then the Mandarin for an evening party.
\end{quote}

\subsection*{You'll want to know: jùhuì}
\emph{jùhuì} — \textquotedblleft{}get-together; party\textquotedblright{}.

\begin{quote}
\emph{zhōumò wǒmen yǒu yí ge jùhuì nǐ lái ma} — We're having a get-together at the weekend: are you coming?
\end{quote}

Its characters come later; for now it is a word for the ear.

\begin{cousinweb}
\emph{jùhuì} is built from \emph{jù} 'to gather' and \zh{会} 'to meet', both inherited from Middle Chinese, \emph{dzjuX} and \emph{hwajH} (Guangyun, in Baxter's transcription); the dictionary glosses the first 'many' and 'to gather', and quotes an older book that glosses it with the second.
\end{cousinweb}

\begin{grammarlens}[title={What you've built}]
One more word for talking about plans.
\end{grammarlens}

\subsection*{Your turn}
\begin{itemize}
\raggedright
  \item \emph{Say it:} \emph{jùhuì}
  \item \emph{Say it:} \emph{jùhuì}, once more
  \item \emph{Say it:} \emph{zhōumò wǒmen yǒu yí ge jùhuì nǐ lái ma} — We're having a get-together at the weekend: are you coming?
  \item \emph{Recall:} say the Mandarin for to send regards to, then the Mandarin for an evening party, then say \emph{jùhuì} again
\end{itemize}

\begin{tcolorbox}[breakable,colback=teal!4,colframe=teal!35!black,title={Before you move on}]
What does jùhuì mean? (\textquotedblleft{}Get-together; party\textquotedblright{}.) Say it once more.
\end{tcolorbox}
\section[\texorpdfstring{\zh{请客} (qǐngkè)}{请客 (qǐngkè)}]{\zh{请客} (qǐngkè) — to pay for everyone; to treat (the others)}

\label{lesson:ZH-C173-qingke}



\begin{quote}
Before the new one: say the Mandarin for an evening party, then the Mandarin for a get-together.
\end{quote}

\subsection*{You'll want to know: \zh{请客}}
\textbf{\zh{请客}} — \emph{qǐngkè} — \textquotedblleft{}to pay for everyone; to treat (the others)\textquotedblright{}.

\begin{quote}
\textbf{\zh{今天我请客每个人一杯茶}} \emph{jīntiān wǒ qǐngkè měi ge rén yì bēi chá} — It's on me today: a cup of tea each.
\end{quote}

\begin{cousinweb}
\textbf{\zh{请客}} (\emph{qǐngkè}) is built from \zh{请} 'to ask, to invite' and \zh{客} 'guest', both inherited from Middle Chinese, \emph{tshjengX} and \emph{khaek} (Guangyun, in Baxter's transcription): to have the others as your guests, and so to pay.
\end{cousinweb}

\begin{grammarlens}[title={What you've built}]
One more word for talking about plans.
\end{grammarlens}

\subsection*{Your turn}
\begin{itemize}
\raggedright
  \item \emph{Say it:} \emph{qǐngkè}
  \item \emph{Say it:} \emph{qǐngkè}, once more
  \item \emph{Say it:} \emph{jīntiān wǒ qǐngkè měi ge rén yì bēi chá} — It's on me today: a cup of tea each.
  \item \emph{Recall:} say the Mandarin for an evening party, then the Mandarin for a get-together, then say \emph{qǐngkè} again
\end{itemize}

\begin{tcolorbox}[breakable,colback=teal!4,colframe=teal!35!black,title={Before you move on}]
What does \zh{请客} mean? (\textquotedblleft{}To pay for everyone; to treat (the others)\textquotedblright{}.) Say it once more.
\end{tcolorbox}
